% Source: https://artofproblemsolving.com/community/c3h1749458_decreasing_and_increasing_numbers
% Stable identity: aops-topic-1749458-q01
\textbf{(a)} The digit $0$ cannot appear: it would have to be the leading digit. Any set of $n$ distinct digits from $\{1,\dots,9\}$ can be written in increasing order in exactly one way, and this gives a valid $n$-digit number. So the count is $\boxed{\dbinom9n}$ (which is $0$ for $n\ge10$).

\textbf{(b)} Any set of $n$ distinct digits from $\{0,1,\dots,9\}$ written in decreasing order gives exactly one number; since $n\ge2$, the largest digit, which comes first, is not $0$. So the count is $\boxed{\dbinom{10}{n}}$ (for example $\dbinom{10}{2}=45$ two-digit numbers such as $10$ and $98$). (For $n=1$ the convention matters: the nine one-digit numbers $1$--$9$, not $\dbinom{10}1=10$.)
